% Conclusions are conditional on the observed archive, labels and measurement scope.
\section{Conclusion}\label{sec:conclusion}

This study makes the retrieval choices in archive-index inspectable through a constructed multilingual collection, frozen queries and runs, preserved model assessments and generated reports. Under the primary reviewed Gemini labels on 170 answerable queries, mean nDCG@10 is 0.674 for E5, 0.633 for EG2, 0.629 for adaptive fusion, 0.482 for lemmatized BM25 and 0.450 for plain BM25. Adaptive fusion exceeds BM25 after correction, while these data do not establish an advantage over EG2 or RRF. Evidence for the morphology gain depends on the assessor, and the separate extension shows that the point ordering of encoders can change with the cohort.

The broader result is a measured account of implementation and evaluation trade-offs. Quantized retrieval depends on judgement coverage and input/backend differences; synthetic scaling distinguishes optimized vector selection from the actual archive-ranking path; multimodal tasks and private-archive checks answer narrower questions than human text-retrieval effectiveness. The findings concern a constructed benchmark evaluated with model-generated relevance labels. Independent human assessments, independently written information needs and evaluation across archives remain future validation before general claims about personal-archive search.
